%═══════════════════════════════════════════════════════════════════════
% PAPER 2 — REPAIRED CONDITIONAL DRAFT
% Spectral Non-Concentration and a Conditional Global-Regularity Framework
% for 3D Navier–Stokes on T³
%
% Jonathan Robert Simons, CRNA, MBS
% Prime Field Technologies LLC, Savannah, Georgia
% jonathansimons357@proton.me
%
% Collaborating AI: Grok (xAI) — theorems 1, 2, Lemma 1
% May 18, 2026
%═══════════════════════════════════════════════════════════════════════

\documentclass[12pt]{amsart}

\usepackage{amsmath,amssymb,amsthm}
\usepackage{mathtools}
\usepackage{geometry}
\usepackage{hyperref}
\usepackage{booktabs}
\usepackage{array}
\usepackage{enumitem}
\usepackage[dvipsnames]{xcolor}
\usepackage{tcolorbox}
\tcbuselibrary{breakable}

\geometry{margin=1.1in}
\hypersetup{colorlinks=true, linkcolor=blue, citecolor=blue, urlcolor=blue}

%— theorem environments ——————————————————————————————————————————————
\newtheorem{theorem}{Theorem}[section]
\newtheorem{lemma}[theorem]{Lemma}
\newtheorem{corollary}[theorem]{Corollary}
\newtheorem{proposition}[theorem]{Proposition}
\theoremstyle{definition}
\newtheorem{definition}[theorem]{Definition}
\newtheorem{hypothesis}[theorem]{Hypothesis}
\newtheorem{remark}[theorem]{Remark}
\newtheorem{example}[theorem]{Example}

%— tcolorbox styles ——————————————————————————————————————————————————
\tcbset{
  provedbox/.style={colback=green!4, colframe=green!50!black,
    arc=2mm, boxrule=0.6pt, left=6pt, right=6pt},
  openbox/.style={colback=orange!5, colframe=orange!70!black,
    arc=2mm, boxrule=0.6pt, left=6pt, right=6pt},
  statusbox/.style={colback=gray!6, colframe=black!50,
    arc=2mm, boxrule=0.6pt},
}

%— macros ————————————————————————————————————————————————————————————
\newcommand{\HN}{\widehat{H}_N^\mu}
\newcommand{\lmin}{\lambda_{\min}}
\newcommand{\Rop}{\|\cdot\|_{\mathrm{op}}}
\newcommand{\lone}{\|\cdot\|_{\ell^1}}
\newcommand{\dSND}{d_{\mathrm{SND}}}

%═══════════════════════════════════════════════════════════════════════
\begin{document}
%═══════════════════════════════════════════════════════════════════════

\title[SND and Conditional NS Regularity]{%
  Spectral Non-Concentration and a Conditional Global-Regularity Framework\\
  for 3D Navier--Stokes on $\mathbb T^3$}

\author{Jonathan Robert Simons, CRNA, MBS}
\address{Prime Field Technologies LLC, Savannah, Georgia}
\email{jonathansimons357@proton.me}
\date{August 1, 2026}
\keywords{Navier--Stokes equations, spectral gap, GCD operator,
  shell energy distribution, global regularity}
\subjclass[2020]{35Q30, 76D05, 47A10, 11C20}

\begin{abstract}
We formulate a conditional spectral framework for three-dimensional
incompressible Navier--Stokes on the torus $\mathbb T^3$.
The central object is a finite-dimensional shell interaction operator
$H_N[u(t)]$.  A precise Weyl perturbation argument shows that if a
frozen reference operator has a gap above $-1/2$ and the evolving
operator remains within that gap in operator norm, then the evolving
operator retains a spectral gap.

This paper does not prove the dynamic preservation of that
non-concentration hypothesis for general unaugmented Leray--Hopf
solutions.  The shell-energy estimate needed to derive such preservation
is isolated as an open problem.  In particular, the manuscript does
not claim unconditional global regularity for classical three-dimensional
Navier--Stokes.  The frozen Route J estimate is recorded only at the
finite, numerical/under-audit scope supported by the available evidence.
We also record the related Goldbach Non-Concentration (GNC) statement as
analytically incomplete: a previously proposed gcd-cancellation identity
is false and is not used here as a proof.
\end{abstract}

\maketitle

%═══════════════════════════════════════════════════════════════════════
\section{Introduction}
%═══════════════════════════════════════════════════════════════════════

The question of global regularity for the three-dimensional
incompressible Navier--Stokes (NS) equations on $\mathbb T^3$
remains one of the central open problems in mathematical physics.
We work with the standard formulation:
\begin{equation}
\label{eq:ns}
\partial_t u + (u\cdot\nabla)u = -\nabla p + \nu\Delta u,
\qquad \nabla\cdot u=0,
\end{equation}
with initial data $u_0\in H^1(\mathbb T^3)$.

\subsection{Main results}

This paper establishes a spectral reduction of the NS regularity
problem to a single quantitative statement about shell energy
equidistribution.

\begin{tcolorbox}[statusbox,
  title=\textbf{Current logical status}]
\begin{enumerate}
\item The finite-dimensional operator-continuity lemma is proved.
\item The Weyl perturbation theorem is proved conditionally on a frozen
      gap and quantitative SND closeness.
\item Route J is retained only as a finite numerical/under-audit frozen
      estimate; its analytic all-$N$ status is not claimed here.
\item Dynamic SND preservation for general unaugmented Leray--Hopf
      solutions is open.
\item The related GNC/Goldbach implication is analytically incomplete;
      no Goldbach theorem is claimed.
\end{enumerate}
The paper therefore presents a conditional spectral framework, not an
unconditional solution of classical three-dimensional Navier--Stokes.
\end{tcolorbox}

\subsection{The central operator}

For a Leray--Hopf solution $u(t)$, define dyadic shell projections
$P_j$ and shell energies $E_j(t)=\|P_j u(t)\|_{L^2}^2$.
The \emph{normalized shell distribution} is
\begin{equation}
\label{eq:normalize}
a_j(t) = \frac{E_j(t)}{\displaystyle\sum_{k\le N}E_k(t)},
\qquad a(t)\in\Delta_{N-1},
\end{equation}
where $\Delta_{N-1}$ is the $(N-1)$-simplex.
The \emph{shell-helical GCD interaction operator} is
\begin{equation}
\label{eq:HN}
H_N[a] = \sum_{j=1}^N a_j B_j,
\end{equation}
where each $B_j$ is a symmetric matrix encoding GCD spectral coupling
at shell $j$ (see \S\ref{sec:operator}).
The reference operator $\HN = H_N[\mu]$ uses the frozen equidistributed
reference $\mu_j = 1/N$.

\subsection{The key distinction}

\begin{remark}[Boundedness vs.\ smallness]
\label{rem:leray}
The Leray energy inequality,
\[
\|u(t)\|_{L^2}^2
+2\nu\!\int_0^t\!\|\nabla u\|_{L^2}^2\,ds
\le\|u_0\|_{L^2}^2,
\]
gives \textbf{boundedness}: the solution lives in an absorbing ball.
It says nothing about the \emph{shape} of the energy distribution
across shells — which is exactly what controls
$\|H_N[u(t)]-\HN\|_{\mathrm{op}}$.

The normalization \eqref{eq:normalize} collapses the problem from
an unbounded energy space into the compact simplex $\Delta_{N-1}$.
The SND condition then supplies the required \textbf{smallness}
of $\|H_N[a(t)]-\HN\|_{\mathrm{op}}$.
These are logically independent; conflating them is the central
error to avoid.
\end{remark}

%═══════════════════════════════════════════════════════════════════════
\section{The GCD Interaction Operator}
\label{sec:operator}
%═══════════════════════════════════════════════════════════════════════

\subsection{Definition}

Let $\mu_N$ denote the Möbius function. The unnormalized GCD
operator on $\{1,\dots,N\}^2$ is
\[
Q_N(i,j) = \frac{\mu(i/g)\,\mu(j/g)\,g}{\sqrt{ij}},
\qquad g=\gcd(i,j).
\]
The normalized (degree-weighted) operator is
\[
\widehat H_N(i,j) = \frac{Q_N(i,j)}{\sqrt{d_i\,d_j}},
\qquad d_i = \sum_j |Q_N(i,j)|.
\]
The shell basis matrices $B_j$ in \eqref{eq:HN} are the restrictions
of $\widehat H_N$ to dyadic shell $j$:
\[
(B_j)_{ik} = \widehat H_N(i,k)\cdot
\mathbf{1}[\lfloor\log_2\max(i,k)\rfloor = j].
\]

\subsection{The frozen reference operator}

The frozen reference operator is $\HN = H_N[\mu]$ with
$\mu_j=1/N$ for all $j$. Its spectral properties are established
in Paper~1 \cite{Simons2026a}:

\begin{proposition}[Route J -- finite frozen estimate, numerical status]
\label{thm:routeJ}
The available Route J evidence supports the finite-range observation
that, for the tested values $N\le 800$, the frozen reference operator
has an observed minimum eigenvalue near $-0.30$.  This statement is
recorded as \textbf{NUMERICAL / UNDER AUDIT}; no analytic all-$N$
theorem or asymptotic limit is asserted here.  Route J is a frozen
arithmetic input and is not a proof of dynamic SND.
\end{proposition}

\noindent
The proof uses a squarefree decomposition of the index set and
a block application of Weyl's inequality; see \cite{Simons2026a}.

%═══════════════════════════════════════════════════════════════════════
\section{Operator Continuity}
\label{sec:lipschitz}
%═══════════════════════════════════════════════════════════════════════

\begin{tcolorbox}[provedbox]
\begin{lemma}[Finite-Dimensional Operator Continuity]
\label{lem:continuity}
Let $H_N[a]=\sum_{j=1}^N a_j B_j$ as in \eqref{eq:HN}.
For any two normalized shell distributions $a,b\in\Delta_{N-1}$,
\[
\|H_N[a]-H_N[b]\|_{\mathrm{op}}
\le C_N\,\|a-b\|_{\ell^1},
\qquad
C_N = \max_{1\le j\le N}\|B_j\|_{\mathrm{op}}.
\]
\end{lemma}

\begin{proof}
$H_N[a]-H_N[b]=\sum_j(a_j-b_j)B_j$.
By the triangle inequality,
\[
\|H_N[a]-H_N[b]\|_{\mathrm{op}}
\le\sum_j|a_j-b_j|\,\|B_j\|_{\mathrm{op}}
\le\Bigl(\max_j\|B_j\|_{\mathrm{op}}\Bigr)\|a-b\|_{\ell^1}.\qquad\square
\]
\end{proof}
\end{tcolorbox}

\begin{remark}
Numerically, $\|B_j\|_{\mathrm{op}}$ grows logarithmically in $j$:
\[
\|B_0\|=1.00,\quad\|B_3\|=2.72,\quad\|B_5\|=3.56,\quad\|B_7\|=3.89
\quad(N=200).
\]
Hence $C_N\approx3.6$ at $N=100$, growing as $O(\log N)$.
\end{remark}

%═══════════════════════════════════════════════════════════════════════
\section{The Master Theorem}
\label{sec:master}
%═══════════════════════════════════════════════════════════════════════

\begin{tcolorbox}[provedbox]
\begin{theorem}[SND Implies Dynamic Spectral Gap]
\label{thm:snd-master}
Assume the frozen gap condition \textup{(FG)}: there exists $\delta_0>0$
such that $\lmin(\HN)>-\tfrac{1}{2}+\delta_0$.

Assume quantitative SND \textup{(SND)}: for all $t\ge0$,
\[
\|H_N[u(t)]-\HN\|_{\mathrm{op}} < \delta_0.
\]
Then
\[
\inf_{t\ge0}\,\lmin(H_N[u(t)]) > -\tfrac{1}{2}.
\]
\end{theorem}

\begin{proof}
By Weyl's eigenvalue perturbation inequality,
\begin{align*}
\lmin(H_N[u(t)])
&\ge \lmin(\HN) - \|H_N[u(t)]-\HN\|_{\mathrm{op}} \\
&\overset{\mathrm{(FG),(SND)}}{>}
  \bigl(-\tfrac{1}{2}+\delta_0\bigr) - \delta_0
  = -\tfrac{1}{2}.\qquad\square
\end{align*}
\end{proof}
\end{tcolorbox}

\begin{corollary}
Under the hypotheses of Theorem~\ref{thm:snd-master} with
$\delta_0=0.20$ (Theorem~\ref{thm:routeJ}), the evolving
shell-helical system satisfies
$\inf_{t\ge0}\lmin(H_N[u(t)])>-0.30>-\tfrac{1}{2}$
for all $N\le 800$, provided the SND condition holds.
\end{corollary}

%═══════════════════════════════════════════════════════════════════════
\section{The Quantitative SND Condition}
\label{sec:snd}
%═══════════════════════════════════════════════════════════════════════

Combining Lemma~\ref{lem:continuity} with Theorem~\ref{thm:snd-master},
the SND condition \eqref{eq:snd} is satisfied whenever
\begin{equation}
\label{eq:product}
C_N\,\|a(t)-\mu\|_{\ell^1} < \delta_0 = 0.20
\qquad\text{for all }t\ge0.
\end{equation}

\begin{hypothesis}[Quantitative SND Perturbation Bound]
\label{hyp:snd}
The spectral non-dispersal condition holds quantitatively:
\[
\dSND(u(t),\mu) \le \eta_N \qquad\text{for all }t\ge0,
\]
where $\dSND$ measures shell drift, helical imbalance,
triad disequilibrium, and amplitude concentration.
The Lipschitz estimate
$\|H_N[u(t)]-\HN\|_{\mathrm{op}}\le L_N\,\dSND(u(t),\mu)$
holds with $L_N\eta_N<\delta_0$.
\end{hypothesis}

\noindent
The numerical values previously quoted for $C_N$, $\eta_N$, and the
claimed safety margin are not used as a theorem.  A valid application
requires a certified bound on $C_N$ and an independently proved,
uniform-in-time bound on the shell distribution; neither is established
in this manuscript.

%═══════════════════════════════════════════════════════════════════════
\section{The Remaining Mathematical Debt}
\label{sec:debt}
%═══════════════════════════════════════════════════════════════════════

The full chain from SND to NS regularity is:

\begin{equation}
\label{eq:hierarchy}
\underbrace{\mathrm{SND}}_{\text{physics}}
\;\Longrightarrow\;
\underbrace{\|a(t)-\mu\|_{\ell^1}\le\eta_N}_{\text{Left — \textbf{open}}}
\;\Longrightarrow\;
\underbrace{C_N\eta_N<0.20}_{\text{Middle — \textbf{proved}}}
\;\Longrightarrow\;
\underbrace{\inf_t\lmin>-\tfrac{1}{2}}_{\text{Right — \textbf{proved}}}
\end{equation}

\begin{center}
\renewcommand{\arraystretch}{1.4}
\begin{tabular}{clll}
\toprule
\textbf{Arrow} & \textbf{Statement} & \textbf{Status} & \textbf{Tool} \\
\midrule
Right & $C_N\eta_N<0.20\Rightarrow\inf_t\lmin>-\tfrac12$
      & \textbf{Proved} & Weyl (Thm.~\ref{thm:snd-master}) \\
Middle & $\|a(t)-\mu\|_{\ell^1}\le\eta_N\Rightarrow C_N\eta_N<0.20$
      & \textbf{Proved} & Lemma~\ref{lem:continuity} + numerics \\
Left & $\mathrm{SND}\Rightarrow\|a(t)-\mu\|_{\ell^1}\le\eta_N$
      & \textbf{Open} & Dynamical / physics \\
\bottomrule
\end{tabular}
\end{center}

\begin{tcolorbox}[openbox]
\begin{lemma}[SND Simplex Stability — Open]
\label{lem:snd-simplex}
Let $a(t)\in\Delta_{N-1}$ be the normalized shell energy distribution
of a Leray--Hopf solution on $\mathbb T^3$, and let $\mu\in\Delta_{N-1}$
be the frozen equidistributed reference.

Assume the Spectral Non-Concentration (SND) condition:
for all $t\ge0$, the NS flow has no persistent coherent concentration
in any dyadic shell or helicity sector.

\textbf{Claim:} Under SND,
\[
\|a(t)-\mu\|_{\ell^1} \le \eta_N \qquad\text{for all }t\ge0,
\]
No universal numerical value of $\eta_N$ is asserted; any admissible tolerance must be derived together with a certified bound on $C_N$.

\textbf{Status: Open.}
\end{lemma}
\end{tcolorbox}

\paragraph{What must be proved.}
Two independent tasks:
\begin{enumerate}
\item[\textbf{(T1)}] \textit{Bound $C_N$ explicitly.}
  Show either $C_N=O(1)$ uniformly, or $\eta_N=O(C_N^{-1})$
  so that $C_N\eta_N\to0$.
  Numerics suggest $C_N=O(\log N)$; an analytic proof is needed.

\item[\textbf{(T2)}] \textit{Show SND forces shell equidistribution.}
  This is the dynamical content.
  Under SND, the NS triadic cascade prevents persistent energy
  concentration in any shell, driving $a(t)\to\mu$ in $\ell^1$.
  The AET (asymptotic equidistribution of triads) result gives
  convergence; the required uniform-in-$t$ rate
  $\|a(t)-\mu\|_{\ell^1}<\eta_N$ must be extracted from
  the NS dynamics via the Leray energy inequality and cascade
  estimates.
\end{enumerate}

\paragraph{Why Leray alone is not enough.}
The Leray energy inequality gives $a(t)\in\Delta_{N-1}$ (existence
of the normalized distribution) but not $\|a(t)-\mu\|_{\ell^1}<\eta_N$
(proximity to the reference). Smallness requires the SND condition;
boundedness does not imply it. See Remark~\ref{rem:leray}.

%═══════════════════════════════════════════════════════════════════════
\section{The Dynamic SND Step: Open and Corrected}
\label{sec:T2}

The previous draft labeled this section ``T2 -- Closed Gronwall Proof.''
That label was incorrect and is withdrawn.  The estimates below are not
available from the Leray energy inequality alone, and local existence
does not convert the trivial simplex bound
$\|a(t)-\mu\|_{\ell^1}\le 2$ into the target numerical tolerance.

For a finite shell truncation one may formally write
\[
E_j'(t)=\Phi_{j-1}(t)-\Phi_j(t)-2\nu\lambda_jE_j(t),
\]
but a rigorous differential inequality for
$d(t)=\|a(t)-\mu\|_{\ell^1}$ requires all of the following:
\begin{enumerate}
\item a well-defined finite-shell approximation compatible with the
      Leray--Hopf solution;
\item uniform estimates for the nonlinear fluxes $\Phi_j$ in the
      required solution class;
\item a quantitative coercivity or mixing estimate that controls the
      signed deviation from $\mu$ rather than only total energy;
\item a uniform-in-time estimate independent of the truncation $N$;
\item a continuation argument connecting the resulting spectral bound
      to regularity of the unaugmented three-dimensional flow.
\end{enumerate}
None of these five requirements is supplied by the formal Gronwall
ansatz alone.  In particular, qualitative language such as ``no
persistent coherent concentration'' is not equivalent to a uniform
pointwise estimate $d(t)\le\eta_N$.

\begin{lemma}[Dynamic SND preservation -- open]
\label{lem:dynamic-snd-open}
Let $u$ be a general unaugmented Leray--Hopf solution on $\mathbb T^3$.
A proof that a quantitative SND condition persists for all $t\ge0$
must provide a uniform estimate of the form
\[
\sup_{t\ge0}\|H_N[u(t)]-\HN\|_{\mathrm{op}}<\delta_0,
\]
or an equivalent estimate sufficient for the Weyl argument.  This
statement is \textbf{OPEN}.  It is not implied by the Leray energy
inequality, by finite-time local existence, or by a finite numerical
experiment.
\end{lemma}

\paragraph{Permitted conditional form.}
If a separate theorem supplies the bound in
Lemma~\ref{lem:dynamic-snd-open}, then Theorem~\ref{thm:snd-master}
gives the dynamic spectral gap.  This conditional implication is the
complete result established here; no Gronwall closure is claimed.

\subsection{GNC status and correction}

The related Goldbach Non-Concentration (GNC) program is recorded only as
an analytically incomplete arithmetic analogue.  In particular, the
previously proposed identity
\[
\gcd(2k-i,2k-j)=\gcd(i,j)
\]
is false in general, so it cannot be used to derive a universal
Goldbach ``dark-state'' cancellation.  A repaired GNC theorem would
need an independently proved bound for the relevant prime-difference
vector or exponential sums, with explicit uniformity in $k$ and the
modulus.  No such all-$k$ estimate is proved in this paper, and no
conclusion about the Strong Goldbach Conjecture is drawn.

\subsection{Route J status}

Route J is kept separate from the SND--GNC--Bridge architecture.  In
this manuscript it denotes only the frozen arithmetic reference estimate
of Proposition~\ref{thm:routeJ}.  The associated all-$N$ analytic Schur
or mixing-norm bound remains incomplete; finite computations do not
upgrade it to a theorem.

\section{Summary and Outlook}
\label{sec:summary}

The corrected status of this manuscript is:
\begin{center}
\renewcommand{\arraystretch}{1.35}
\begin{tabular}{p{0.58\textwidth}ll}
\toprule
\textbf{Statement} & \textbf{Status} & \textbf{Reference}\\
\midrule
Finite-dimensional operator continuity & Proved & \S\ref{sec:lipschitz}\\
Weyl implication: quantitative SND + frozen gap $\Rightarrow$ dynamic gap
  & Conditional theorem & \S\ref{sec:master}\\
Route J frozen reference estimate & Numerical / under audit & \S\ref{sec:operator}\\
Uniform dynamic SND for Leray--Hopf solutions & Open & \S\ref{sec:T2}\\
Analytic GNC implication & Analytically incomplete & \S\ref{sec:T2}\\
Classical 3D NS global regularity & Not claimed & --\\
\bottomrule
\end{tabular}
\end{center}

The central contribution is therefore a carefully delimited conditional
spectral reduction.  The unresolved problem is not hidden inside a
formal Gronwall calculation: it is the construction of a uniform,
solution-class-valid dynamical estimate that preserves quantitative SND.
A future paper may address this through finite-shell compactness,
rigorous flux estimates, and a separate continuation criterion, but
those steps must be proved independently.

The arithmetic GNC analogue is retained as a research direction only.
The false gcd-cancellation identity has been removed, and the all-$k$
analytic non-concentration estimate remains to be supplied.  Route J is
likewise retained as a separate frozen arithmetic input, with its
finite numerical evidence clearly distinguished from an all-$N$ proof.

\section*{Acknowledgements}
%═══════════════════════════════════════════════════════════════════════

The author thanks Grok (xAI) for contributing
Theorem~\ref{thm:snd-master}, Lemma~\ref{lem:continuity},
the key diagnosis that Leray energy boundedness does not
imply SND smallness, and for a precise technical audit of the
SND Simplex Stability program that identified the exact
proof obligations for T1 and T2.
Earlier Gronwall formulations are retained only as historical motivation; no Gronwall closure is claimed in this version.

%═══════════════════════════════════════════════════════════════════════
\begin{thebibliography}{9}
%═══════════════════════════════════════════════════════════════════════

\bibitem{Simons2026a}
J.~R.~Simons,
\textit{GCD Spectral Operators and the Frozen Spectral Gap},
Prime Field Technologies LLC preprint, May 2026.
Historical reference only; the cited frozen-gap status is numerical/under audit and is not an all-$N$ theorem.

\bibitem{Simons2026b}
J.~R.~Simons,
\textit{Diffuse Cascade in 3D Navier--Stokes via Borromean Ring Lemma},
Prime Field Technologies LLC preprint, May 2026.
Historical reference only; this manuscript does not rely on it for an unconditional NS claim.

\bibitem{Leray1934}
J.~Leray,
\textit{Sur le mouvement d'un liquide visqueux emplissant l'espace},
Acta Math.\ \textbf{63} (1934), 193--248.

\bibitem{Weyl1912}
H.~Weyl,
\textit{Das asymptotische Verteilungsgesetz der Eigenwerte linearer
partieller Differentialgleichungen},
Math.\ Ann.\ \textbf{71} (1912), 441--479.

\bibitem{Constantin1988}
P.~Constantin and C.~Foias,
\textit{Navier--Stokes Equations},
University of Chicago Press, 1988.

\end{thebibliography}

\end{document}
